\section{Метод решения}
Решение разделено на общую логику лабораторной и платформенный слой. Родительская программа организует ввод команд и обмен данными, дочерняя выполняет проверку чисел и запись результата. Вызовы ОС, связанные с процессами, каналами и выходным файлом, вынесены в библиотеку \texttt{os}.

\subsection{Создание процесса и перенаправление потоков}
Родитель создаёт два неименованных канала. Первый передаёт строки от родителя к ребёнку, второй возвращает однобайтовые статусы. Стандартный ввод ребёнка связывается с читающим концом первого канала, стандартный вывод --- с пишущим концом второго. Диагностические сообщения выводятся в стандартный поток ошибок и не смешиваются со статусами.

\begin{center}
\begin{tabular}{lll}
Направление & Передаваемые данные & Назначение \\
\hline
Родитель $\rightarrow$ ребёнок & Строка с числом & Команда \\
Ребёнок $\rightarrow$ родитель & Один байт & Статус обработки \\
Ребёнок $\rightarrow$ файл & Число и перевод строки & Результат \\
\end{tabular}
\end{center}

После запуска ребёнка родитель закрывает свои копии читающего конца первого канала и пишущего конца второго. Родителю остаются только конец отправки команд и конец чтения ответов. Это необходимо, в частности, для корректного обнаружения закрытия канала~\cite{winpipes,linuxpipes}.

\subsection{Протокол обмена}
Дочерний процесс сначала открывает выходной файл и отправляет статус готовности \texttt{R}. Только после этого родитель начинает передавать числа. Такая последовательность позволяет обнаружить ошибку открытия файла до начала обработки команд.

\begin{center}
\begin{tabular}{cl}
Статус & Значение \\
\hline
\texttt{R} & Дочерний процесс готов к работе \\
\texttt{W} & Составное число записано в файл \\
\texttt{I} & Число 0 или 1 пропущено \\
\texttt{E} & Ошибка формата числа или выход за диапазон \texttt{int} \\
\texttt{S} & Получено простое или отрицательное число \\
\end{tabular}
\end{center}

На каждую переданную строку родитель ожидает один статус. Статус немедленно отправляется ребёнком со сбросом буфера вывода. При \texttt{S} родитель прекращает ввод, закрывает каналы и ожидает завершения дочернего процесса. Если пользователь завершает ввод, закрытие пишущего конца первого канала позволяет ребёнку обнаружить конец потока.

\subsection{Проверка простоты}
После разбора строки отрицательное значение приводит к завершению работы. Для неотрицательного значения проверяется простота. Значение 2 является простым, остальные чётные числа больше 2 --- составными. Для нечётного числа проверяются нечётные делители начиная с 3.

Проверка продолжается, пока выполняется условие \(d \leq n/d\). Оно эквивалентно ограничению \(d^2 \leq n\) для положительных чисел, но не требует умножения, которое могло бы вызвать переполнение. При обнаружении делителя число признаётся составным. Временная сложность проверки составляет \(O(\sqrt{n})\), дополнительная память --- \(O(1)\).

\section{Описание программы}
\subsection{Структура файлов}
Пути приведены относительно корня проекта \texttt{Operation-Systems}.

\begin{itemize}
\item \texttt{lab2/code/parent/main.cpp} --- ввод имени файла и команд, запуск ребёнка, приём статусов и ожидание завершения.
\item \texttt{lab2/code/child/main.cpp} --- обработка команд, проверка чисел, отправка статусов и запись результатов.
\item \texttt{lab2/code/common.cpp} --- общий разбор целого числа и проверка простоты.
\item \texttt{lab2/code/include/common.h} --- объявления функций общей логики.
\item \texttt{lab2/code/include/protocol.h} --- константы статусов обмена.
\item \texttt{lab2/code/include/resources.h} --- управление временем жизни файловых ресурсов.
\item \texttt{common/code/os/os.h} --- платформенно-независимый интерфейс.
\item \texttt{common/code/os/os.cpp} --- полная запись данных и формирование командной строки.
\item \texttt{common/code/os/os\_win.cpp} --- реализация интерфейса для Windows.
\item \texttt{common/code/os/os\_posix.cpp} --- реализация интерфейса для Linux.
\end{itemize}

\subsection{Основные типы и функции}
\texttt{os::FileHandle} и \texttt{os::ProcessHandle} скрывают платформенное представление ресурсов. В реализации Windows используются дескрипторы \texttt{HANDLE}, в Linux --- файловые дескрипторы и идентификаторы процессов. Прикладная логика обращается к ним через функции общего интерфейса.

\texttt{os::ProcessStartInfo} содержит путь программы, аргументы и параметры перенаправления стандартных потоков. \texttt{os::CreateProc} запускает ребёнка и возвращает дескриптор процесса. \texttt{os::WaitProc} ожидает завершения, а \texttt{os::CloseProc} освобождает ресурс процесса.

\texttt{os::CreatePipe} создаёт канал. \texttt{os::PipeRead} читает данные, \texttt{os::PipeWrite} выполняет отдельную запись, \texttt{os::WriteAll} повторяет запись до передачи всех байтов. \texttt{os::OpenWrite} открывает выходной файл на запись. \texttt{os::CloseFile} освобождает файловый ресурс.

Обёртка \texttt{UniqueFile} реализует RAII: при выходе из области видимости вызывает закрытие ресурса. Копирование владельца запрещено, поэтому один ресурс не закрывается двумя независимыми владельцами.

\subsection{Используемые средства Windows}
\begin{itemize}
\item \texttt{CreatePipe}, \texttt{SetHandleInformation} --- создание каналов и управление наследованием.
\item \texttt{DuplicateHandle} --- подготовка наследуемых стандартных дескрипторов.
\item \texttt{InitializeProcThreadAttributeList}, \texttt{UpdateProcThreadAttribute}, \texttt{DeleteProcThreadAttributeList} --- явный список наследуемых дескрипторов.
\item \texttt{CreateProcessW} --- запуск дочерней программы с перенаправлением стандартных потоков.
\item \texttt{ReadFile}, \texttt{WriteFile}, \texttt{CreateFileW} --- ввод-вывод и создание выходного файла.
\item \texttt{WaitForSingleObject}, \texttt{GetExitCodeProcess}, \texttt{CloseHandle} --- ожидание, проверка кода завершения и освобождение ресурсов.
\item \texttt{GetModuleFileNameW}, \texttt{MultiByteToWideChar}, \texttt{WideCharToMultiByte} --- определение пути программы и преобразование кодировки путей.
\end{itemize}

В \texttt{CreateProcessW} передаётся изменяемый буфер командной строки. Путь к ребёнку определяется относительно исполняемого файла родителя. В текущей сборке используются имена \texttt{lab2\_parent.exe} и \texttt{lab2\_child.exe}. Дескриптор первичного потока закрывается сразу после создания процесса; это не останавливает сам поток. Дескриптор процесса сохраняется для ожидания завершения~\cite{winprocess}.

\subsection{Используемые средства Linux}
\begin{itemize}
\item \texttt{pipe2}, \texttt{fcntl} --- создание каналов и настройка \texttt{FD\_CLOEXEC}.
\item \texttt{fork}, \texttt{dup2}, \texttt{execv} --- создание процесса, перенаправление стандартных потоков и запуск программы ребёнка.
\item \texttt{open}, \texttt{read}, \texttt{write}, \texttt{close} --- работа с файлами и каналами.
\item \texttt{waitpid} --- получение результата завершения ребёнка.
\item \texttt{readlink} --- получение пути текущей программы через \texttt{/proc/self/exe}.
\item \texttt{sigaction} --- настройка обработки \texttt{SIGPIPE}.
\end{itemize}

Внутри запуска процесса дополнительно используется временный служебный канал для передачи ошибки подготовки или \texttt{execv} родителю. При успешном \texttt{execv} он закрывается благодаря \texttt{FD\_CLOEXEC}. Этот канал не участвует в обмене пользовательскими числами.

\subsection{Обработка ошибок}
Ошибки системных операций преобразуются в исключения с названием операции и кодом ОС. Внешние обработчики выводят диагностику в \texttt{stderr} и устанавливают ненулевой код завершения. Родитель проверяет не только статусы протокола, но и код завершения ребёнка.

Частичные записи обрабатываются в \texttt{WriteAll}. В Linux чтение и запись повторяются при \texttt{EINTR}; сигнал \texttt{SIGPIPE} игнорируется, чтобы закрытие читающего конца можно было обработать как ошибку записи, а не аварийное завершение процесса~\cite{linuxpipes}.
